\documentclass{report}

\usepackage{amsmath}
\usepackage{amssymb}
\usepackage{booktabs}
\usepackage{adjustbox}

\begin{document}
\chapter{Research Question and Economic Motivation}

\section{Research Question and Economic Motivation}

Financial market volatility plays a central role in financial economics, as it reflects uncertainty surrounding asset prices and influences investment decisions, portfolio allocation, and risk management. Accurate volatility forecasts are particularly important for investors and financial institutions seeking to assess market risk and manage their exposure to unexpected price fluctuations. However, forecasting volatility remains challenging because financial markets are interconnected and may respond simultaneously to changes in economic conditions, commodity prices, and investor expectations.
\vspace{1em}
Commodity markets represent a potentially valuable source of information for understanding fluctuations in equity market volatility. In particular, oil prices are closely related to economic activity, production costs, and inflationary pressures, while gold is frequently associated with financial uncertainty and investors' demand for alternative assets. Natural gas and agricultural commodities, such as corn, may also reflect developments in energy markets, supply conditions, and broader economic uncertainty. Consequently, changes in commodity market volatility may contain information about future fluctuations in equity market volatility that is not fully captured by the historical behaviour of the equity market itself.

\vspace{1em}
Against this background, this study addresses the following research question:

\vspace{1em}

\textbf{“To what extent does incorporating information from commodity futures market volatility improve the forecasting accuracy of S&P 500 futures volatility compared with models relying solely on historical equity market information?”}

\vspace{1em}
To investigate this question, the analysis focuses on the E-mini S\&P 500 futures market and examines whether volatility information from West Texas Intermediate (WTI) crude oil and gold futures improves predictions of future equity market volatility. Natural gas and corn futures are additionally considered to assess whether including a broader range of commodity markets provides further predictive value.
\vspace{1em}
The economic motivation for this research is based on the possibility that volatility is transmitted across financial markets through common economic shocks, changing expectations, and shifts in investor risk perceptions. For example, heightened uncertainty in energy markets may coincide with increased uncertainty regarding corporate profitability, inflation, and economic growth, potentially affecting equity market volatility. Similarly, changes in gold market volatility may reflect broader financial uncertainty. These relationships suggest that models incorporating information from multiple markets could capture relevant volatility dynamics that would otherwise be overlooked.
\vspace{1em}
An important distinction, however, is that relationships between commodity and equity market volatility do not necessarily imply improved forecasting performance. Even when volatility movements across markets are statistically related, commodity information may provide little additional predictive value once historical equity volatility is taken into account. Therefore, the objective of this study is not to establish causal relationships between commodity markets and equity volatility, but rather to evaluate whether cross-market information contributes to more accurate out-of-sample forecasts.
\vspace{1em}
This question is relevant to financial economics because improvements in volatility forecasting can support more informed risk assessment, portfolio management, and financial decision-making. Furthermore, examining the incremental predictive value of commodity volatility provides insight into the extent to which information from interconnected financial markets can enhance conventional forecasting approaches. By comparing models based exclusively on equity market information with models incorporating additional commodity market variables, this study aims to assess whether greater model complexity translates into meaningful improvements in predictive accuracy.
\vspace{1em}
\chapter{Data}
\label{sec:data}

We use daily data on futures contracts from five markets. The E-mini S\&P~500 (ES) represents the equity market, while crude oil (CL), gold (GC), natural gas (NG) and corn (C) represent the commodity markets. We include oil and gold as two of the largest commodity markets, and natural gas and corn to cover energy and agricultural risk.

\vspace{1em}

Realized variances and closing prices come from the Volatility Archive for Realized Estimates (VOLARE; Cipollini et al., 2026). For futures, VOLARE computes its measures from intraday mid-quote prices, the average of the bid and ask quotes supplied by Kibot. Our volatility measure is the 5-minute realized variance, the sum of squared 5-minute returns over each trading day, where each return is based on the last available quote at or before the 5-minute mark. Weekends and major holidays are not included. All realized variances are log-transformed, which compresses the large spikes in their distribution. We also compute daily returns from futures closing prices.

For the bond market, we use the daily market yield on US Treasury securities at 10-year constant maturity (DGS10). The series is published by the Board of Governors of the Federal Reserve System in its H.15 release and is estimated from the yields of Treasury securities across maturities. We obtain it from FRED, the public economic database of the Federal Reserve Bank of St.~Louis. Yields are quoted in percent and are missing on days when the bond market is closed. As no intraday bond data are available, we measure bond volatility by the size of the daily change in the yield, expressed in basis points. Since yields are reported in whole basis points, this change is exactly zero on many days. Because these zero changes are retained, the proxy is not log-transformed. To allow for publication timing, the yield change is delayed by one observed Treasury date, so that each forecast uses only information available before its target date.

The VOLARE data run from 28 September 2009 to 31 August 2026. The early gold series contains missing observations and extreme realized variances that coincide with implausible intraday prices. To avoid this irregular early period, we start in January 2011. We keep only dates on which ES, CL, GC, NG and the Treasury yield are all observed, without filling in missing values. This leaves 3,910 daily observations from 5 January 2011 to 31 August 2026. The first 3,127 observations, up to 12 July 2023, form the initial training sample, and the remaining 783 observations form the test sample on which all models are evaluated. Corn, which enters only the HAR, is missing on one additional date in the training sample, while the test sample is unaffected. Lags refer to the previous common observation. Beyond these steps, we do not remove observations, so crisis episodes such as the April 2020 oil price collapse remain in the sample.

\end{document}